\documentclass[11pt]{article}
\usepackage[T1]{fontenc}
\usepackage[margin=1in]{geometry}
\usepackage{amsmath,amssymb,amsthm,microtype}
\usepackage[colorlinks=true,allcolors=blue]{hyperref}
\newtheorem{theorem}{Theorem}[section]
\newtheorem{lemma}[theorem]{Lemma}
\title{Small Boolean circuits satisfying constraints on their truth tables}
\author{Research manuscript}
\date{22 November 2026}
\begin{document}\maketitle
\begin{abstract}
Partition the $N=2^n$ entries of a Boolean truth table into cells and prescribe the number of ones in every cell. If a fraction $2^{-R}$ of all tables of the required weight meet the quotas, every placement of the cells admits a satisfying table computable with $R/n+O(N\log n/n^2)$ gates. A uniform permutation of the cells attains this bound to leading order when $R$ is a positive fraction of $N$. We count the different realizations available at larger budgets. We also identify the typical circuit size needed to realize a prescribed mean at every address by a distribution of Boolean circuits. When the largest cell has $o(N)$ entries, the leading bound holds for arbitrary relations among the entries of each cell. As a consequence, a decomposable predicate of subexponential size that accepts a fraction $2^{-rN}$ of a weight layer accepts some table of circuit size at most $(r+o(1))N/n$.
\end{abstract}
\section{How large is the smallest witness?}
A Boolean circuit on $n$ address bits outputs one of the $N=2^n$ entries of a word $y$. Write $C(y)$ for its minimum gate count. The upper bounds use binary AND, OR and counted unary NOT; the lower bounds allow all binary Boolean gates and complemented wires. All logarithms have base two. We allow unrestricted fanout. The information cost of a property on a set $U$ of candidate tables is $\log(|U|/|\mathcal F|)$, where $\mathcal F\subseteq U$ comprises the feasible tables.

For a partition into cells of sizes $M_i$, prescribe $a_i$ ones in each cell. Put
\[
 W=\sum_i a_i,\qquad H=\log\binom NW,\qquad
 T=\prod_i\binom{M_i}{a_i},\qquad R=H-\log T.
\]
Here $0\le a_i\le M_i$. There is no restriction on the number of cells. Zero information can be misleading if measured only with limiting frequencies. With cells of even size $m$ and quota $m/2$, the prescribed common frequency is $1/2$, but the exact information per site tends to $1-m^{-1}\log\binom{m}{m/2}$. This positive constant contributes to the minimum circuit size.

\begin{theorem}\label{th:quota}
For every placement and every $0\le\lambda\le1$, at least $2^{\lambda\log T-O(N\log n/n)}$ distinct satisfying words have
\[
 C(y)\le \frac{R+\lambda\log T}{n}+O\!\left(\frac{N\log n}{n^2}\right).
\]
In particular, at $\lambda=0$ at least one such word exists. The constants are absolute. If $R/N\to r>0$, then both a uniform coordinate permutation of the specification, in probability, and a worst placement give minimum size $(r+o(1))N/n$.
\end{theorem}
Write $Z_s$ for the number of feasible words of circuit size at most $s$. At $s\sim cN/n$ with $R/N\to r<c$ and $H/N\to h$, the minimum over placements and the typical value of $N^{-1}\log Z_s$ both converge to $\min(c,h)-r$. The typical statement uses a uniformly random permutation of a fixed specification. Endpoints with $c=r$ require the additive-error formulation.

\section{A dictionary for the next block}
Read the addresses in natural order, in blocks of dyadic length $n/16\le b\le n/8$. At a prefix of whole blocks let $T_p$ be the number of feasible completions and let $T_{py}$ be the count after adding the next block $y$. Uniform conditional completion gives probabilities $p_y=T_{py}/T_p$. Group the possible $y$ by weight $w$ and by $\lfloor-\log p_y\rfloor$. A feasible block has probability at least $N^{-b}$, since each next feasible bit in its cell has conditional probability at least $1/N$. Among at most $(b+1)(bn+1)$ groups choose one of greatest mass, with support $S$ and $K=|S|$. Every $y\in S$ satisfies
\begin{equation}\label{eq:credit}
 \log K\ge\log T_p-\log T_{py}-O(\log n).
\end{equation}
For a given block, the local partition, the remaining sizes and quotas of its touched cells, its weight and bucket specify $S$. Untouched cells cancel in the conditional probability. There are at most $b^b(N+1)^{2b}(b+1)(bn+1)$ descriptions. The logarithm of this count is $O(bn)$, independent of the number of cells.

Fix each weight layer of the block and write $M=\binom bw$. There is a common permutation of its $M$ words such that for every candidate $S$ and every $1\le L\le K$, the first $\min(M,\lceil CbnLM/K\rceil)$ words contain at least $L$ elements of $S$. Indeed, a random prefix that stops before $M$ has expected intersection size at least $CbnL$. The hypergeometric lower tail is $\exp(-\Omega(CbnL))$. Choose an absolute $C$ large enough for the union over descriptions and $L\le2^b$. A full-length prefix always succeeds. Choose $L=\max(1,\lfloor K^\lambda/2\rfloor)$ and retain just $L$ eligible words. Each has an index costing at most
\[
 \log M-(1-\lambda)\log K+O(\log n)\quad\text{bits}.
\]
Along a path, the sum of $\log M$ is at most $H$, since the product of its block weight layers embeds in the global weight layer. Equation~\eqref{eq:credit} telescopes from $T$ to one. Hence the sum of the index widths is at most $R+\lambda\log T+O(N\log n/n)$. If each allowed next word is chosen uniformly, every completed path has probability at most $2^{-\lambda\log T+O(N\log n/n)}$ by the same telescope. Paths are distinct tables, which gives the list count.

\paragraph{From bits to gates.} Each block has a weight, an index width at most $b$, and an index. Group blocks by the first two quantities. For each group and bit of its index, store the bit values as a vector ordered by occurrence of the group. For addressing these vectors use consecutive superblocks of $\Theta(n^8)$ blocks. Each block needs $O(\log n)$ directory bits for its group and its ordinal within that group in its superblock. For each superblock and each group record the count of earlier occurrences, using $O(n)$ bits. The aggregate directory has $O(N\log n/n)$ bits. This design charges the different index widths by their sum. A single lookup on a common block address padded to the maximum width would lose the leading constant.

There are polynomially many vectors. Split each length into at most $n+1$ aligned dyadic pieces. For an $m$-bit address, Lupanov lookup uses $(2^m/m)(1+O(\log m/m))$ gates. For completeness, split off $\lceil3\log m\rceil$ low address bits, precompute minterms on both sides and group the low minterms into batches of $\lfloor m-8\log m\rfloor$. Precompute every subset OR in each batch, then select the prescribed batch outputs for each high-address minterm and take their union. The leading cost is the number of input positions divided by the batch size. Direct lookup handles bounded $m$. Fix $A$ sufficiently large compared with the degree of the number of pieces. All pieces of length $<N/n^A$ cost $O(N/n^2)$ in total by direct lookup. The others have $m=n-O(\log n)$. The directory, multiplexers, and a common truth table of size $\operatorname{poly}(n)2^b$ decoding the block permutation fit the stated error. This proves the circuit upper bound.

\paragraph{The lower bound.} A size-$s$ binary circuit has a depth-first backwards description with at most $s+1$ pointers to previously named gates, input variables, or constants. Other fields cost $O(s)$ bits. At $s=O(N/n)$ the number of different tables is thus at most $2^{sn+O(N\log n/n)}$. Every particular table of weight $W$ is feasible under a uniform coordinate permutation with probability exactly $T/\binom NW$. The first moment gives both the typical lower bound on size and the typical upper bound on the list length. The maximum possible number of distinct words in the layer supplies the other bound used in $\min(c,h)$.

\section{The cost of prescribing every coordinate mean}\label{sec:marginal}
Let $p=(p_1,\ldots,p_N)\in[0,1]^N$. We seek a distribution on small deterministic address circuits whose output at each site $i$ has expectation exactly $p_i$. Equivalently, $p$ must lie in the convex hull of their Boolean truth tables. Let $C_{\rm mix}(p)$ be the least gate budget allowed for every member of the mixture. Write $H_2$ for binary entropy and let
\[
 J(p)=N H_2(\bar p)-\sum_{i=1}^{N}H_2(p_i),\qquad
 \bar p=N^{-1}\sum_i p_i .
\]
The quantity $J(p)/N$ is the mutual information between a uniform site and its prescribed random output. A constant vector has $J(p)=0$ and uses a mixture of the two constant circuits. In contrast, if every $p_i$ equals zero or one, any exact mixture is supported on the single word $p$.

\begin{theorem}\label{th:marginal}
For every $p\in[0,1]^N$,
\[
 C_{\rm mix}(p)\le J(p)/n+O(N\log n/n^2).
\]
For a uniform permutation $\pi$ of its coordinates, $n C_{\rm mix}(p\circ\pi)/N-J(p)/N$ tends to zero in probability, uniformly over sequences $p$. The same asymptotic holds for a most expensive permutation.
\end{theorem}
\begin{proof}
Write $s$ for the claimed upper bound with a sufficiently large absolute constant. By separation of a point and a compact convex polytope, it suffices, for every real score vector $a$, to build one $y$ of circuit size at most $s$ with $\sum_i a_i y_i\ge\sum_i a_i p_i$. For a positive score define the favorable outcome to be one and put $p_i'=p_i$; at a negative score use zero and put $p_i'=1-p_i$. Either choice works at a zero score. Apart from the same additive constant on both sides, the scores now sum $|a_i|$ times the favorable indicator or probability.

Put $B=n^2$ and round each $p_i'$ upwards to $q_i=\lceil Bp_i'\rceil/B$. Form strata by $q_i=q$ and the favorable bit, and sort the sites within each stratum by decreasing $|a_i|$. Divide the stratum into consecutive blocks of length $B$, except possibly its last block. Pin the first and last blocks to the favorable bit. In each other block prescribe exactly $q B$ favorable outcomes. Every resulting word has the following prefix property in each stratum: its number of favorable outcomes on the prefix is at least the sum of the original favorable probabilities on that prefix. Indeed the initial full block creates slack $(1-q)B$ against the constant fractional benchmark $q$. In a later full block the largest possible deficit on a partial prefix is at most $q(B-qB)\le(1-q)B$, and the complete block has zero deficit. The pinned last block creates no deficit from its starting value. A single shorter stratum is pinned entirely. Summation by parts with the decreasing weights now gives the required score inequality.

These are exact quotas of the form in Theorem~\ref{th:quota}. At most $O(B^2)$ sites are pinned because there are $2(B+1)$ strata. The mean output weight differs from $\bar p$ by $O(B^{-1}+B^2/N)$. In the nonpinned blocks, $\log\binom B{qB}=B H_2(q)+O(\log(B+1))$. The binary entropy continuity bound $|H_2(u)-H_2(v)|\le H_2(|u-v|)$ and $H_2(p)=H_2(1-p)$ give
\[
 H-\log T\le J(p)+O(N\log n/n^2+\operatorname{poly}(n)).
\]
The quota theorem and the separation criterion prove exact containment. In particular no uniform lower bound on $p_i$ or $1-p_i$ was required.

For the matching typical lower bound, bin the original coordinates by rounding $p_i$ to a grid of mesh $1/n^2$. Write $Z$ for the bin of a uniformly chosen position, and permute these bins along with their original values. For a fixed Boolean word $y$, write $Y$ for its output at that same uniform position. At every empirical joint type, the probability of the type under this permutation is at most $2^{-NI(Z;Y)+O(n^3)}$. In fact it equals a product of cell binomial coefficients divided by the overall weight-layer size. There are at most $(N+1)^{O(n^2)}$ possible types. A union bound over all circuits of size at most $s=O(N/n)$ and all types shows with probability $1-o(1)$ that every such $y$ has $NI(Z;Y)\le sn+o(N)$; the error is uniform and deterministic in $s$ by also union-bounding its $O(N)$ possible values.

For fixed positions, mixing words amounts to mixing their binary-output channels with the input-bin distribution held fixed. Mutual information is convex in this channel, by the log-sum inequality. A mixture with means $p\circ\pi$ has channel output probability within $1/n^2$ of every prescribed $p_i$ in each bin. Entropy continuity therefore gives $NI(Z;Y)=J(p)+o(N)$ for that mixture. Hence $sn\ge J(p)-o(N)$ with the asserted probability. The deterministic upper bound then also identifies the worst placement.
\end{proof}

\paragraph{Pointwise approximation of a Boolean target.} Fix $0<D<1/2$. Every Boolean target $f$ has a distribution on circuits of size at most $(1-H_2(D))N/n+O(N\log n/n^2)$ with error exactly $D$ at every address: apply the theorem to $p_i=D+(1-2D)f_i$. Conversely, if a distribution on circuits of size $s$ achieves at most $D$ error at every address, some member is within Hamming distance $DN$ of $f$ by averaging. The volume of a Hamming ball and the circuit census show that almost every $f$ needs $s\ge(1-H_2(D)-o(1))N/n$. Thus this size is asymptotically optimal even under the one-sided relaxation of the numerical error bound from equality to inequality. The model permits an arbitrary distribution over gate lists. Requiring a single fixed gate list with explicit coin wires is a different model: at error less than one half each such list determines only one target by pointwise majority.

\paragraph{Exactly uniform choice among a prescribed number of circuits.}
Suppose $m p_i$ is an integer at every site for some positive integer $m$. We can require the mixture to choose uniformly from $m$ circuits while retaining the leading upper bound whenever $m\to\infty$. More quantitatively, put $\xi=m^{-1/2}+1/n$. For $\xi<1/2$ the additional size over $J(p)/n$ is $O((H_2(\xi)+\log n/n)N/n)$, uniformly in $m$.

To prove this, cut the addresses into between $n$ and $2n$ equal dyadic chunks. On each chunk independently consider $m$ draws from the exact-mean ensemble of Theorem~\ref{th:marginal}, applied at the shorter address length. The expected sum over chunk sites of the difference in absolute value between the actual and desired number of ones is at most the chunk length times $\sqrt m/2$, by binomial variance. Fix draws satisfying this aggregate bound on each chunk. Flip exactly that many entries among the $m$ rows to correct every site count. We next permute the rows of each whole chunk, preserving the counts at its sites. At each chunk assign its rows in decreasing order of the number of corrections to global rows in increasing order of their current correction totals. The range of these totals stays at most the maximum chunk length: for two previously ordered totals, opposing the order of their new increments bounds the new absolute difference by the larger of the two old ranges. Thus every global row has at most $N/(2\sqrt m)+N/n$ corrections.

On each row concatenate the originally drawn chunk circuits. Each has address width $n-O(\log n)$ and their total information budgets are at most $J(p)$ by concavity of $H_2$. The address selectors and total lookup errors fit $O(N\log n/n^2)$ gates. The correction table on the row is sparse. Apply Theorem~\ref{th:quota} with singleton quotas to that particular correction word to synthesize it with at most $H_2(\xi)N/n+O(N\log n/n^2)$ gates. XOR it with the concatenated circuit. The $m$ corrected rows have exactly the prescribed coordinate means. Repeated rows may be retained as distinct tickets in the uniform choice. For an arbitrary real $p$, rounding to the $1/m$ grid gives coordinate error at most $1/m$ and increases $J(p)$ by at most $O(NH_2(1/m))$.

\section{Arbitrary predicates inside small cells}
Let each cell impose an arbitrary Boolean relation among its entries. We may additionally prescribe the exact total weight $W$. Write $T$ for the number of compatible words on that layer, again with $H=\log\binom NW$.
\begin{theorem}\label{th:relations}
Suppose the largest cell has $o(N)$ entries and $T>0$. For every $0\le\lambda\le1$ there are at least $2^{\lambda\log T-o(N)}$ compatible words with $C(y)\le [H-(1-\lambda)\log T+o(N)]/n$. The losses are uniform given an upper bound tending to zero on the largest cell's fraction of $N$.
\end{theorem}
\begin{proof}
The preceding dictionary only requires a short description of every possible next-block probability bucket. We construct such descriptions even when the predicates themselves are expensive to evaluate.

\paragraph{Spreading the cells.} Let $p$ be the maximum cell fraction. In the array of $N/b$ blocks by $b$ positions, independently translate the high address bits in each position-column by a random XOR shift. Each cell's occupancy of a fixed physical block is a sum of independent indicators of mean at most $pb$. Put $t=\lceil10 n/\log(1/p)\rceil=o(n)$. For large $n$ its upper-tail probability beyond $t$ is at most $(epb/t)^t\le2^{-3n}$. A union bound fixes shifts so every block meets every cell in at most $t$ sites. Computing this permutation or its inverse needs only polynomially many gates in $n$.

\paragraph{Recording global weight.} Put $k=\lfloor\sqrt n\rfloor$. Fix the weights of all cells larger than $k$ to some profile retaining at least $T'=T2^{-o(N)}$ feasible words. Indeed, $\sum_{M_i>k}\log(M_i+1)=O(N\log(k+1)/k)$. The remaining global constraint on the smaller cells is their total residual weight. At a prefix the uniform feasible completion now factors into independent large-cell laws and one law on the small cells jointly.

For each large cell touched in the current block, record the conditional probabilities of its at most $2^t$ local patterns by rounding their negative logarithms upwards to multiples of $1/b$, with a sentinel above $2b$. Any block pattern hitting a sentinel has probability at most $2^{-2b}$. Summing over all block patterns loses at most $2^{-b}$. For the others the products of the rounded probabilities have error at most a factor of two.

For a touched small cell record the set of conditionally permitted completions on its at most $k$ remaining sites, and which sites occur in the current block. Altogether these touched small cells have at most $bk$ remaining sites. Let $W'$ be the residual required weight in all small cells. For each $0\le j\le bk$, record the integer part of the logarithm of the number of compatible assignments of weight $W'-j$ on the untouched small cells, or a zero sentinel. Each integer needs $O(n)$ bits. For a candidate block word, sum over compatible patterns on the rest of the touched small cells, weighting each by the recorded approximation for the needed untouched weight. The resulting mass is within a factor of two of the true numerator, and normalizing distorts probabilities by at most another constant factor. The entire local description, including the arrangement of sites, uses $D=\operatorname{poly}(n)(2^t+2^k)$ bits. This counts possible descriptions. A circuit for a satisfying word uses only its short index.

\paragraph{Using the dictionary.} Discard all additional block words whose estimated probability is less than $2^{-b-C}$ for a sufficiently large constant $C$. At least half the mass survives. Group remaining words by weight and dyadic estimated probability; a heaviest bucket has actual mass $\Omega(1/b^2)$ and its atom probabilities differ by only a fixed factor. Thus (\ref{eq:credit}) still holds. The logarithm of the number of candidate buckets is $O(D+\log n)$. Oversampling by $O(D+n)$ in the hypergeometric argument gives a valid common dictionary with an additional $O(\log(D+n))=O(t+k+\log n)=o(n)$ bits per block. The path count and packed circuit lookup of Section~2 give the theorem.
\end{proof}
The size condition is needed for arbitrary predicates. On a dyadic address subcube of some fixed density $\alpha>0$ exclude all restrictions realizable with at most $\alpha N/(3n)$ gates and leave other sites free. A uniform balanced table matches any specific subcube pattern with probability at most $O(\sqrt N)2^{-\alpha N}$, by conditioning independent bits on total weight. The circuit census implies that the excluded fraction is exponentially small, yet every remaining table needs more than $\alpha N/(3n)$ gates. Quotas retain the short conditional descriptions of Section~2 even in this regime.

\section{Prescribing the joint law inside each cell}
Let the cells be disjoint with $m_i$ sites, $\sum_i m_i=N$, and let $p_i$ be an arbitrary distribution on binary words of length $m_i$. A mixture of address circuits must have precisely marginal law $p_i$ on cell $i$. Different cells may be correlated in the mixture. Let $\rho=N^{-1}\sum_i\mathbb E_{p_i}|Y_i|$ and define
\[
 J_{\rm cell}(p)=N H_2(\rho)-\sum_i H(p_i).
\]
\begin{theorem}\label{th:cellmix}
If $\max_i m_i=o(N)$, the prescribed cell laws are realized exactly by a distribution on Boolean circuits, each of size at most $(J_{\rm cell}(p)+o(N))/n$.
\end{theorem}
\begin{proof}
A separating linear score on the space of marginal vectors has the form $\sum_i a_i(Y_i)$, for arbitrary real-valued functions $a_i$ on their cells. We construct a small table scoring at least $\sum_i\mathbb E_{p_i}a_i$. Put $m_0=\max_i m_i$ and group whole cells into clusters of size comparable to $L=\lceil\sqrt{Nm_0}\rceil$. A greedy filling to size at least $L$, adding the final short remainder to the preceding cluster, gives all cluster sizes between $L$ and $3L$ for large $n$.

Inside one cluster choose independently from $p_i'=(1-\eta)p_i+\eta\delta_{z_i}$, where $z_i$ maximizes $a_i$ and $\eta=(m_0/L)^{1/8}$. Define the nonnegative score loss $d_i=\max a_i-a_i$. The expected aggregate loss under this perturbation is $(1-\eta)$ times its original expectation. Markov's inequality shows that the cluster meets its original expected score with probability at least $\eta$; zero original loss gives probability one. The bound is independent of the units or ranges of the scores.

For the same draw put $S_i=-\log p_i'(Y_i)$. Its tail obeys $\Pr[S_i>m_i+u]\le2^{-u}$ for $u\ge0$, by counting at most $2^{m_i}$ atoms. Thus the variance of its sum is $O(\sum_i m_i^2)=O(m_0L)$, as is the variance of the total number of ones in the cluster. The expected surprise is at least $(1-\eta)\sum_i H(p_i)$ by concavity, and the expected weight differs from its original value by at most $3\eta L$. Chebyshev's inequality with tolerance $\eta L$ fails with probability $O((m_0/L)^{3/4})=o(\eta)$ for either quantity. Thus with probability at least $\eta/2$ a word meets the score, has its original expected weight up to $O(\eta L)$, and has surprise at least $\sum_i H(p_i)-O(\eta L)$. There are at least $2^{\sum_i H(p_i)-o(L)}$ such cluster words. Restricting them to a single weight loses at most a factor $3L+1$.

Use these words as an allowed relation on each cluster. Their product has $T\ge2^{\sum_i H(p_i)-o(N)}$ words, all of weight $N(\rho+o(1))$. Every word in the product attains the separating score. The clusters have $o(N)$ sites, so Theorem~\ref{th:relations} at $\lambda=0$ supplies the required small word. Separation proves exact simultaneous feasibility of the marginal vectors.
\end{proof}


\begin{theorem}[Exact asymptotic rate for arbitrary local marginal laws]\label{th:qconv}
For each length separately, define
\[
 R_*(p)=\min\{r:p\in\operatorname{conv}\{q:J_{\rm cell}(q)\le r\}\}.
\]
Here $q$ ranges over the product of the local probability simplices for the actual cells. Let $s_{\min}$ be the smallest budget for a distribution on Boolean circuits with marginals $p$. Every placement has $s_{\min}\le(R_*(p)+o(N))/n$ if $\max_i m_i=o(N)$. Under a uniform permutation of the sites, $n s_{\min}/N-R_*(p)/N$ tends to zero in probability. There is no restriction on the number of distinct marginal laws or on their atom probabilities.
\end{theorem}
\begin{proof}
The sublevel sets are compact. In each finite dimension, take a finite convex decomposition of $p$ into laws $q$ of cost at most $R_*(p)$. The error in Theorem~\ref{th:cellmix} is uniform in $q$. Mixing its circuit ensembles gives the upper bound.

For the converse take $r<R_*(p)$. A separating functional gives a real score $A(y)=\sum_i a_i(y_i)$ whose prescribed expectation $A(p)$ exceeds the maximum for marginal vectors in the convex hull at level $r$. For each integer weight $W$, let $E_W$ be the set of weight-$W$ words with score at least $A(p)$. If it is nonempty, take the uniform distribution on $E_W$ and let $q_i$ be its cell marginals. Then $A(q)\ge A(p)$, so $J_{\rm cell}(q)>r$. Entropy subadditivity gives
\[
 \log |E_W|\le\sum_i H(q_i)<N H_2(W/N)-r.
\]
It follows that $|E_W|/\binom NW\le(N+1)2^{-r}$, uniformly in $W$. When the layout is permuted each fixed Boolean word of weight $W$ lands uniformly on this layer relative to the score. A union bound over all words of circuit size $s$, using the census, shows that with probability at least $1-(N+1)2^{sn-r+o(N)}$ none obtains the prescribed expected score. Then no distribution over them can realize the marginals. Choose $r$ between $sn$ and $R_*(p)$ with a positive linear gap whenever $sn\le R_*(p)-\epsilon N$. This proves the claim, including varying cell sizes and target laws.
\end{proof}

For singleton cells $J_{\rm cell}$ is mutual information times $N$, convex in the binary-output channel. Its sublevel already is convex and the theorem reduces to Section~\ref{sec:marginal}. For general cells, the distinction matters. Put $u=N^{-1}\sum_i\mathbb E_{q_i}|Y_i|$. Then
\[
 J_{\rm cell}(q)=\sum_i D_2\!\left(q_i\,\Vert\,\operatorname{Ber}(u)^{\otimes m_i}\right),
\]
where $D_2$ is relative entropy in bits. Its zero set consists of product-bit cell laws sharing one bias $u$. The exact finite-length condition $R_*(p)=0$ therefore means there is one measure $\pi$ on $[0,1]$ for which, simultaneously in every cell,
\[
             p_i=\int \operatorname{Ber}(u)^{\otimes m_i}\,d\pi(u).
\]
A common random choice can create dependence inside the prescribed cell marginals.

\section{Witnesses for decomposable predicates}
A decomposable negation normal form (DNNF) has literal leaves and AND and OR gates. Each AND combines disjoint variable sets. The predicate now receives all $N$ entries as independent input variables. Its number of wires measures the complexity of \emph{testing} a property, whereas the circuit for a witness computes an entry from its $n$ address bits.
\begin{theorem}
Let a DNNF of size $2^{o(N)}$ have $T>0$ satisfying words on a weight layer of logarithmic size $H$. The list and circuit bound of Theorem~\ref{th:relations} hold for this predicate. In particular, a predicate of this size that certifies complexity greater than $cN/n$ accepts at most a fraction $2^{-cN+o(N)}$ of any weight layer. The theorem also applies to syntactically read-once nondeterministic branching programs of size $2^{o(N)}$.
\end{theorem}
\begin{proof}
Binarize the circuit and smooth it by completing smaller OR-input scopes with tautologies, adding missing variables at the root. Polynomial overhead keeps the resulting size $S$ subexponential; take $S\ge N$. In an accepting trace, an OR keeps the chosen child's scope, while an AND splits into disjoint scopes. Starting from a scope of size $m>2$, follow an accepting child at an OR and a larger child at an AND. At the first node with at most $2m/3$ variables, its size is at least $m/3$. For each possible such gate $g$, its satisfying set times its possible contexts lies inside the root's satisfying set. These products cover the root. This is the balanced rectangle-cover argument of Bova et al.\ \cite{b}.

To repeat the split, the context must itself have a small DNNF. Along a path from the output to $g$, conjoin the off-path siblings at AND gates. The OR over all paths of these conjunctions has a circuit with $O(S)$ wires. In reverse topological order, a child context takes the OR over its parents' contexts, conjoining the other sibling's gate for an AND parent. By induction its scope is exactly the root scope minus the child scope; each new AND is therefore decomposable and each OR is smooth. The sibling gates retain their original representations. Every satisfying assignment of such a context extends any satisfying assignment of $g$ along the corresponding path.

Choose $B=\lceil\sqrt{N\log(S+2)}\rceil=o(N)$ and recursively split any factor larger than $B$, using a fixed rule to break ties. A lineage is split $O(\log(N/B))$ times. Each terminal scope has at least $B/3$ entries, so each branch of the cover uses $O(N/B)$ splits. Even allowing a constant factor growth in circuit size per level, the number of products is at most
\[
  2^{O((N/B)(\log S+\log(N/B)))}=2^{o(N)}.
\]
Some product retains at least $T2^{-o(N)}$ satisfying words in the layer, and all its individual factors have scope at most $B$. Apply Theorem~\ref{th:relations}. A read-once nondeterministic program converts to a DNNF using OR for guesses and guards by a literal for each tested branch. Syntactic read-once ensures the guarded successor does not reuse the tested variable.
\end{proof}
For an unrestricted weight, choose a weight layer containing at least $T/(N+1)$ satisfying words and use $H\le N$. For any fixed predicate or cell system, uniformly permuting all its input coordinates preserves its representation size and the count $T$. The first-moment bound from Section~2 then gives the corresponding typical lower bound when the information cost is linear in $N$.

\section{Related work}
Pippenger derived an information-theoretic circuit bound for approximate realization of partially specified Boolean functions. Andreev studied the entropy of a nondeterministic finite function with allowed multi-bit outputs associated to each row. Andreev, Clementi and Rolim treated approximation of total and partial Boolean functions, and Chashkin studied the complexity of partial functions. In the cell problem a circuit receives the original bit address. A map from that address to its cell and its position within the cell can itself require Shannon-order complexity. The dictionary proof above avoids evaluating that map. The density is measured inside the required weight layer. Exact channel simulation with common randomness also relates simulation cost to mutual information. Harsha, Jain, McAllester and Radhakrishnan measure expected communication; here every sampled address circuit has a gate budget \cite{hjmr}. The balanced-cover theorem for decomposable circuits used in the preceding section is due to Bova, Capelli, Mengel and Slivovsky. We do not claim a definitive priority determination for all variants.\par
\begin{thebibliography}{9}
\bibitem{p} N. Pippenger, \href{https://doi.org/10.1007/BF01683269}{Information theory and the complexity of Boolean functions}, Mathematical Systems Theory 10, 129--167.
\bibitem{a} A. E. Andreev, \href{https://doi.org/10.7146/brics.v1i2.21668}{Complexity of nondeterministic functions}, BRICS Report Series, 1994.
\bibitem{acr} A. E. Andreev, A. E. F. Clementi and J. D. P. Rolim, \href{https://eccc.weizmann.ac.il/eccc-reports/1995/TR95-041/index.html}{Optimal bounds for the approximation of Boolean functions and some applications}, Theoretical Computer Science, 1997.
\bibitem{hjmr} P. Harsha, R. Jain, D. McAllester and J. Radhakrishnan, \href{https://www.tcs.tifr.res.in/~prahladh/papers/HJMR/HJMR2010.pdf}{The communication complexity of correlation}, IEEE Transactions on Information Theory, 2010.
\bibitem{cs} I. Csisz\'ar, \href{https://doi.org/10.1109/18.720546}{The method of types}, IEEE Transactions on Information Theory, 1998.
\bibitem{ch} A. Chashkin, \href{https://doi.org/10.20948/mvk-2024-152}{On circuit complexity of partial Boolean functions}, Mathematical Problems of Cybernetics, 2024.
\bibitem{b} S. Bova, F. Capelli, S. Mengel and F. Slivovsky, \href{https://ijcai.org/proceedings/2016/147}{Knowledge compilation meets communication complexity}, IJCAI 2016.
\end{thebibliography}
\end{document}
